% ============================================================
% Section 59: 离子晶体肖特基缺陷的平衡浓度与体积变化
% ============================================================
\newpage
\section{固体理论：离子晶体肖特基缺陷的平衡浓度与体积变化}
\label{sec:schottky-defect}

\begin{modelbox}[题目：离子晶体肖特基缺陷的统计热力学]
(1) 在离子晶体中，由于电中性的要求，肖特基缺陷都是成对地产生。令 $n$ 代表正负离子空位的对数，$u$ 是产生一对肖特基缺陷所需要的能量，$N$ 为整个离子晶体中正负离子对的数目。试证明：
\begin{equation}
n = N e^{-u/2k_{\!B}T}.
\end{equation}

(2) 肖特基缺陷是晶体内部的原子（离子）跑到晶体表面上，而使原来的位置变为空位缺陷，因此肖特基缺陷将引起晶体体积增大。试由上述结论估计离子晶体中有肖特基缺陷时其体积相对变化 $\Delta V/V$；

(3) 设 NaCl 只有肖特基缺陷，在 $800^{\circ}\mathrm{C}$ 时用 X 射线衍射测定 NaCl 的离子间隙，由此确定的质量密度计算得的分子量为 $58.430$，而用化学方法测定得分子量为 $58.454$。试计算在 $800^{\circ}\mathrm{C}$ 时缺陷的相对浓度以及正负离子空位对应的形成能。
\end{modelbox}

\begin{tipbox}[考点快速分析]
\begin{itemize}
    \item \textbf{电中性约束}：离子晶体中，单独的正离子空位或负离子空位都会破坏电中性，因此肖特基缺陷必须以\textbf{正负离子空位对}的形式出现
    \item \textbf{与弗仑克尔缺陷的对比}：两者都涉及两个独立的选择过程（弗仑克尔：格点+间隙；肖特基：正离子格点+负离子格点），故指数上均有 $1/2$
    \item \textbf{体积效应}：肖特基缺陷的物理图像是内部离子``迁移到表面''，晶体整体增加一个原胞体积，这是区别于弗仑克尔缺陷的关键特征
    \item \textbf{实验测定}：X 射线衍射测晶格常数（得到``表观密度''），化学法测真实质量与体积（得到``真实密度''），两者差异反映空位导致的体积膨胀
\end{itemize}
\end{tipbox}

\begin{derivationbox}[标准解题过程]
\textbf{Step 1: 肖特基缺陷平衡浓度的证明}

\textit{自由能变化。}肖特基缺陷是热平衡缺陷。在恒温恒压下，系统的自由能 $F = U - TS$ 在平衡时取极小值。产生 $n$ 对正负离子空位时：
\begin{itemize}
    \item 内能增加：$\Delta U = nu$
    \item 熵增加：来源于组态熵
\end{itemize}

\textit{组态熵计算。}正离子空位和负离子空位各自独立地在 $N$ 个格点上分布。正离子空位的排列方式数为
\begin{equation}
W_+ = \binom{N}{n} = \frac{N!}{(N-n)!\,n!}.
\end{equation}
负离子空位的排列方式数相同：$W_- = W_+$。总微观状态数为
\begin{equation}
W = W_+ W_- = \left[\frac{N!}{(N-n)!\,n!}\right]^2.
\label{eq:schottky-W}
\end{equation}
组态熵为
\begin{equation}
\Delta S = k_{\!B}\ln W = 2k_{\!B}\ln\!\left[\frac{N!}{(N-n)!\,n!}\right].
\end{equation}

\textit{平衡条件。}自由能变化
\begin{equation}
\Delta F = \Delta U - T\Delta S = nu - 2k_{\!B}T\ln\!\left[\frac{N!}{(N-n)!\,n!}\right].
\end{equation}
利用斯特林近似 $\ln x! \approx x\ln x - x$，对 $n$ 求导并令其为零：
\begin{equation}
\frac{\partial\Delta F}{\partial n} = u - 2k_{\!B}T\ln\!\left(\frac{N-n}{n}\right) = 0.
\end{equation}
解得
\begin{equation}
\frac{N-n}{n} = e^{-u/2k_{\!B}T}.
\end{equation}
在缺陷浓度很低时（$n\ll N$），$N-n\approx N$，故
\begin{equation}
\boxed{n = N e^{-u/2k_{\!B}T}.}
\label{eq:schottky-approx}
\end{equation}
\textit{证毕。}

\textbf{Step 2: 肖特基缺陷引起的体积相对变化}

肖特基缺陷的形成机制是内部离子迁移到表面，在原位留下空位。每产生一对正负离子空位，相当于把一个原胞体积的离子对从内部搬到了表面，晶体体积增加一个分子体积 $v$。

设无缺陷时晶体体积 $V_0 = Nv$，产生 $n$ 对空位后体积增加 $\Delta V = nv$。因此相对体积变化为
\begin{equation}
\frac{\Delta V}{V_0} = \frac{nv}{Nv} = \frac{n}{N}.
\end{equation}
代入 \eqref{eq:schottky-approx} 得
\begin{equation}
\boxed{\frac{\Delta V}{V_0} = e^{-u/2k_{\!B}T}.}
\end{equation}
\textit{结论}：肖特基缺陷引起的体积相对膨胀\textbf{直接等于缺陷浓度}。

\textbf{Step 3: 实验数据求缺陷浓度与形成能}

\textit{基本原理。}X 射线衍射测量晶格常数，得到的是实际原子间距。由此计算出的密度对应\textbf{实际晶体体积}（含空位膨胀后的体积 $V$）。化学方法测量宏观密度（质量/体积），但计算分子量时需使用真实体积。若误用实际体积，则得到偏小的表观分子量。

设真实分子量为 $m$（无缺陷时的理论值，化学法测得），表观分子量为 $m'$（X 射线数据计算值）。由密度公式
\begin{equation}
\rho = \frac{M}{V} = \frac{m'N_{\!A}}{v_{\mathrm{cell}}},\qquad
\rho_0 = \frac{mN_{\!A}}{v_{\mathrm{cell}}},
\end{equation}
得
\begin{equation}
\frac{m}{m'} = \frac{\rho_0}{\rho} = \frac{V}{V_0}.
\end{equation}
即真实分子量与表观分子量之比等于晶体实际体积与理想体积之比。

\textit{计算缺陷浓度。}已知
\begin{equation}
\frac{m}{m'} = \frac{58.454}{58.430} \approx 1.00041,
\end{equation}
因此
\begin{equation}
\frac{V}{V_0} = 1 + \frac{\Delta V}{V_0} = 1.00041
\quad\Longrightarrow\quad
\frac{\Delta V}{V_0} = 0.00041 = 4.1\times 10^{-4}.
\end{equation}
由 Step 2 的结论，缺陷相对浓度即为
\begin{equation}
\boxed{\frac{n}{N} = \frac{\Delta V}{V_0} = 4.1\times 10^{-4}.}
\end{equation}

\textit{计算形成能 $u$。}由 $n/N = e^{-u/2k_{\!B}T}$，取对数：
\begin{equation}
u = 2k_{\!B}T\ln\!\left(\frac{N}{n}\right).
\end{equation}
代入数值：
\begin{align}
T &= 800^{\circ}\mathrm{C} = 1073\,\mathrm{K},\\
k_{\!B} &= 8.617\times 10^{-5}\,\mathrm{eV/K},\\
\ln(N/n) &= \ln(1/0.00041) = \ln 2439 \approx 7.80,
\end{align}
得
\begin{equation}
u = 2\times 8.617\times 10^{-5}\times 1073\times 7.80
\approx \boxed{1.44\,\mathrm{eV}}.
\end{equation}
\textit{注}：NaCl 肖特基缺陷形成能的实验值约为 $2.0\sim 2.5\,\mathrm{eV}$。本题数据简化处理得到的 $1.44\,\mathrm{eV}$ 为估算数量级。若按教材原始数据 $k_{\!B}=8.615\times 10^{-5}\,\mathrm{eV/K}$ 并取 $\ln 2439\approx 7.799$，可得 $u\approx 1.44\,\mathrm{eV}$。
\end{derivationbox}

\begin{formulabox}[答案汇总]
\begin{center}
\begin{tabular}{ll}
\toprule
\textbf{问题} & \textbf{答案} \\
\midrule
(1) 平衡浓度 & $n = N e^{-u/2k_{\!B}T}$ \\
(2) 体积相对变化 & $\displaystyle\frac{\Delta V}{V_0} = \frac{n}{N} = e^{-u/2k_{\!B}T}$ \\
(3) $800^{\circ}\mathrm{C}$ 缺陷浓度 & $n/N = 4.1\times 10^{-4}$ \\
(3) 形成能 & $u \approx 1.44\,\mathrm{eV}$ \\
\bottomrule
\end{tabular}
\end{center}
\end{formulabox}

\begin{insightbox}[物理图像：为什么肖特基缺陷会增大体积？]
弗仑克尔缺陷是内部原子在晶体内部``搬家''（格点$\to$间隙），晶体总体积\textbf{不变}——原子没离开晶体，只是换了位置。

肖特基缺陷的本质不同：内部原子\textbf{跑到了晶体表面}，在原位留下空位。这意味着：
\begin{enumerate}
    \item 晶体内部少了 $n$ 个原胞的内容物
    \item 但晶体表面增加了 $n$ 个原子层，使整体外轮廓扩大
    \item 净效果：体积增加 $n$ 个分子体积 $v$
\end{enumerate}
\textbf{定量关系}：$\Delta V/V_0 = n/N$。这个简洁结果说明，通过测量晶格常数（X射线）和宏观密度（化学法）的差异，可以直接读出缺陷浓度——这是固体物理中经典的\textbf{点缺陷实验诊断方法}。
\end{insightbox}

\begin{thinkbox}[考场速记：缺陷统计问题的对比表]
\begin{center}
\begin{tabular}{lccc}
\toprule
\textbf{缺陷类型} & \textbf{微观图像} & \textbf{微观状态数} & \textbf{平衡浓度} \\
\midrule
弗仑克尔 & 格点$\to$间隙（内部搬家） & $\displaystyle\binom{N}{n}\binom{N'}{n}$ & $\displaystyle n=\sqrt{NN'}e^{-u/2k_{\!B}T}$ \\
肖特基（离子晶体） & 内部$\to$表面（留下空位对） & $\displaystyle\binom{N}{n}^2$ & $\displaystyle n=N e^{-u/2k_{\!B}T}$ \\
肖特基（单质） & 内部$\to$表面（留下单空位） & $\displaystyle\binom{N}{n}$ & $\displaystyle n=N e^{-u/k_{\!B}T}$ \\
\bottomrule
\end{tabular}
\end{center}
\vspace{4pt}
\textbf{记忆口诀}：
\begin{itemize}
    \item 两种子系统 $W=W_1W_2$（弗仑克尔或离子肖特基）$\Rightarrow$ 指数带 $1/2$
    \item 一种子系统 $W=W_1$（单质肖特基）$\Rightarrow$ 指数不带 $1/2$
    \item 体积变化只与肖特基有关：$\Delta V/V_0 = n/N$
\end{itemize}
\end{thinkbox}
